\section{The prediction ledger}
\label{sec:ledger}

Table~\ref{tab:ledger} counts every scored prediction by phase. Of 59 verdicts, 40
were confirmed, 4 were mixed, 14 were not confirmed, and 1 could not be tested with
our instruments. We group the misses below by what caused them; the reports discuss
each one, including those that fit no group, such as the order of TTFT across
budgets.

\begin{table}[t]
  \centering
  \caption{Scored predictions by phase. Phase B scores ten signatures for each of two
    workload regimes. Milestone 5's C8 predictions are counted in Part~2. Where a
    report gives a split or qualified verdict (Phase D predictions 2 and 7), we count
    it as mixed.}
  \label{tab:ledger}
  \begin{tabular}{lrrrrr}
    \toprule
    Phase & Scored & Confirmed & Mixed & Not confirmed & Not testable \\
    \midrule
    B: workload regimes          & 20 & 14 & 0 & 5 & 1 \\
    C: concurrency crossover     & 5  & 3  & 0 & 2 & 0 \\
    C: GPU contention            & 7  & 4  & 2 & 1 & 0 \\
    D: sweep and token budget    & 11 & 6  & 2 & 3 & 0 \\
    Milestone 5, C32 block       & 7  & 6  & 0 & 1 & 0 \\
    Milestone 5, Part 2          & 9  & 7  & 0 & 2 & 0 \\
    \midrule
    Total                        & 59 & 40 & 4 & 14 & 1 \\
    \bottomrule
  \end{tabular}
  % src: ledger built from each report's prediction table; verdict words as stated in each report, except Phase D 2 ("met as specified, uninformative") and 7 ("not confirmed for p50; confirmed for p99"), counted as mixed
\end{table}

\paragraph{The hypothesis was wrong.}
We predicted that the prefill-heavy workload would be dominated by prefill, which a
cost estimate refutes (Section~\ref{sec:load}), and that time before the first token
would dominate past saturation, when the added work went into fuller steps instead.
The step-size account predicted that prompts shorter than the budget would make steps
smaller; the scheduler packed them, at both C32 and C8. The budget cap itself behaved
as predicted.

\paragraph{The instrument could not see the effect.}
vLLM's iteration-token histogram could not show per-step chunk sizes, so Phase D's
manipulation check failed as specified and its step-size test was inconclusive. That
miss led to the per-step log. \texttt{nvidia-smi} utilization is a busy fraction
averaged over the driver's own sample period \citep{nvidiasmi}, so no polling rate
could resolve 40\,ms prefills; we did not replace it. The step log's interval then
turned out to be confounded with the previous step (Section~\ref{sec:stepcost}), and
Part~2's execution time corrects this at C8 and C1 only.

\paragraph{The scheduler and the load generator shaped the result.}
Two queue predictions missed: queue growth past saturation was sub-linear (not
confirmed), and contention raised queue time less than prefill time (mixed). Both fit
vLLM's policy of scheduling running requests first and filling the rest of the budget
with waiting prompts, which turns waiting into longer steps. The closed-loop client
also caps how many requests can wait.

\paragraph{A test was too strict for its data.}
The rank correlation between step size and step interval cleared its threshold at C32
(0.87) but not at C8 (0.57). Forward and reverse condition orders differed by
0.1--2.5\% in latency and throughput, enough to fail a strict range-overlap test in
the first workload phase. With three runs per order, two ranges fail to overlap by
chance one time in ten for each metric even without an order effect, so across about
ten metrics a few failures were expected.

Scoring every prediction changed the study. The per-step crossover and the context
measurements were tested directly only because an earlier prediction failed or
prompted a post hoc reading.
% src: ledger section (lessons as stated in each report); Phase B stable-signature row; 2 x 3!3!/6! = 0.1
